\numoptstrue
\yearpaper{2020}{13 September 2020 \textperiodcentered{} Test Booklet Code E1 \textperiodcentered{} Physics Q 136--180 \textperiodcentered{} 180 Marks}{In this paper Physics is the last section; the original question numbers are kept.}
\setcounter{qno}{135}
\begin{multicols}{2}

\QQ{For which one of the following, Bohr model is \textbf{not} valid?}
{\opt{Hydrogen atom}{Singly ionised helium atom (He$^+$)}{Deuteron atom}{Singly ionised neon atom (Ne$^+$)}}

\QQ{The ratio of contributions made by the electric field and magnetic field components to the intensity of an electromagnetic wave is: ($c$ = speed of electromagnetic waves)}
{\opt{$c:1$}{1:1}{$1:c$}{$1:c^2$}}

\QQ{The Brewster's angle $i_b$ for an interface should be:}
{\opt{$0^\circ<i_b<30^\circ$}{$30^\circ<i_b<45^\circ$}{$45^\circ<i_b<90^\circ$}{$i_b = 90^\circ$}}

\QQ{A cylinder contains hydrogen gas at pressure of $249\unit{kPa}$ and temperature $27^\circ$C. Its density is: ($R = 8.3\unit{J\,mol^{-1}K^{-1}}$)}
{\opt{$0.5\unit{kg/m^3}$}{$0.2\unit{kg/m^3}$}{$0.1\unit{kg/m^3}$}{$0.02\unit{kg/m^3}$}}

\QQ{A ray is incident at an angle of incidence $i$ on one surface of a small angle prism (with angle of prism $A$) and emerges normally from the opposite surface. If the refractive index of the material of the prism is $\mu$, then the angle of incidence is nearly equal to:}
{\opt{$\dfrac A{2\mu}$}{$\dfrac{2A}\mu$}{$\mu A$}{$\dfrac{\mu A}2$}}

\QQ{Two cylinders A and B of equal capacity are connected to each other via a stop cock. A contains an ideal gas at standard temperature and pressure. B is completely evacuated. The entire system is thermally insulated. The stop cock is suddenly opened. The process is:}
{\opt{isothermal}{adiabatic}{isochoric}{isobaric}}

\QQ{The energy equivalent of $0.5\unit{g}$ of a substance is:}
{\opt{$4.5\times10^{16}\unit{J}$}{$4.5\times10^{13}\unit{J}$}{$1.5\times10^{13}\unit{J}$}{$0.5\times10^{13}\unit{J}$}}

\QQ{A body weighs $72\unit{N}$ on the surface of the earth. What is the gravitational force on it, at a height equal to half the radius of the earth?}
{\opt{48\,N}{32\,N}{30\,N}{24\,N}}

\QQ{The solids which have the negative temperature coefficient of resistance are:}
{\opt{metals}{insulators only}{semiconductors only}{insulators and semiconductors}}

\QQ{The phase difference between displacement and acceleration of a particle in a simple harmonic motion is:}
{\opt{$\pi$ rad}{$\dfrac{3\pi}2$ rad}{$\dfrac\pi2$ rad}{zero}}

\QQ{A screw gauge has least count of $0.01\unit{mm}$ and there are 50 divisions in its circular scale. The pitch of the screw gauge is:}
{\opt{0.01\,mm}{0.25\,mm}{0.5\,mm}{1.0\,mm}}

\QQ{In a guitar, two strings A and B made of same material are slightly out of tune and produce beats of frequency 6\,Hz. When tension in B is slightly decreased, the beat frequency increases to 7\,Hz. If the frequency of A is 530\,Hz, the original frequency of B will be:}
{\opt{523\,Hz}{524\,Hz}{536\,Hz}{537\,Hz}}

\QQ{Two particles of mass $5\unit{kg}$ and $10\unit{kg}$ respectively are attached to the two ends of a rigid rod of length $1\unit{m}$ with negligible mass. The centre of mass of the system from the $5\unit{kg}$ particle is nearly at a distance of:}
{\opt{33\,cm}{50\,cm}{67\,cm}{80\,cm}}

\QQ{Find the torque about the origin when a force of $3\hat j\unit{N}$ acts on a particle whose position vector is $2\hat k\unit{m}$.}
{\opt{$6\hat i\unit{N\,m}$}{$6\hat j\unit{N\,m}$}{$-6\hat i\unit{N\,m}$}{$6\hat k\unit{N\,m}$}}

\QQ{Light with an average flux of $20\unit{W/cm^2}$ falls on a non-reflecting surface at normal incidence having surface area $20\unit{cm^2}$. The energy received by the surface during time span of 1 minute is:}
{\opt{$10\times10^3\unit{J}$}{$12\times10^3\unit{J}$}{$24\times10^3\unit{J}$}{$48\times10^3\unit{J}$}}

\QQ{A spherical conductor of radius $10\unit{cm}$ has a charge of $3.2\times10^{-7}\unit{C}$ distributed uniformly. What is the magnitude of electric field at a point $15\unit{cm}$ from the centre of the sphere? $\left(\dfrac1{4\pi\varepsilon_0} = 9\times10^9\unit{N\,m^2/C^2}\right)$}
{\opt{$1.28\times10^4\unit{N/C}$}{$1.28\times10^5\unit{N/C}$}{$1.28\times10^6\unit{N/C}$}{$1.28\times10^7\unit{N/C}$}}

\QQ{In a certain region of space with volume $0.2\unit{m^3}$, the electric potential is found to be $5\unit{V}$ throughout. The magnitude of electric field in this region is:}
{\opt{zero}{0.5\,N/C}{1\,N/C}{5\,N/C}}

\QQ{The increase in the width of the depletion region in a p-n junction diode is due to:}
{\opt{forward bias only}{reverse bias only}{both forward bias and reverse bias}{increase in forward current}}

\QQ{A $40\unit{\mu F}$ capacitor is connected to a $200\unit{V}$, $50\unit{Hz}$ ac supply. The rms value of the current in the circuit is, nearly:}
{\opt{1.7\,A}{2.05\,A}{2.5\,A}{25.1\,A}}

\QQ{The mean free path for a gas, with molecular diameter $d$ and number density $n$ can be expressed as:}
{\opt{$\dfrac1{\sqrt2\,n\pi d}$}{$\dfrac1{\sqrt2\,n\pi d^2}$}{$\dfrac1{\sqrt2\,n^2\pi d^2}$}{$\dfrac1{\sqrt2\,n^2\pi^2d^2}$}}

\QQ{For transistor action, which of the following statements is \textbf{correct}?}
{\opt{Base, emitter and collector regions should have same doping concentrations.}{Base, emitter and collector regions should have same size.}{Both emitter junction as well as the collector junction are forward biased.}{The base region must be very thin and lightly doped.}}

\QQ{Light of frequency 1.5 times the threshold frequency is incident on a photosensitive material. What will be the photoelectric current if the frequency is halved and intensity is doubled?}
{\opt{doubled}{four times}{one-fourth}{zero}}

\QQ{When a uranium isotope ${}^{235}_{\ 92}$U is bombarded with a neutron, it generates ${}^{89}_{36}$Kr, three neutrons and:}
{\opt{${}^{144}_{\ 56}$Ba}{${}^{91}_{40}$Zr}{${}^{101}_{\ 36}$Kr}{${}^{103}_{\ 36}$Kr}}

\QQ{The energy required to break one bond in DNA is $10^{-20}\unit{J}$. This value in eV is nearly:}
{\opt{6}{0.6}{0.06}{0.006}}

\QQ{Two bodies of mass $4\unit{kg}$ and $6\unit{kg}$ are tied to the ends of a massless string. The string passes over a pulley which is frictionless (see figure). The acceleration of the system in terms of acceleration due to gravity ($g$) is:
\fig{\begin{tikzpicture}[scale=0.7]
\draw (0,0) circle (0.5); \fill (0,0) circle (1pt); \draw (-0.5,0) -- (-0.5,-1.2); \draw (0.5,0) -- (0.5,-1.9);
\draw (-0.8,-1.2) rectangle (-0.2,-1.7); \node[lbl,left] at (-0.8,-1.45) {4\,kg};
\draw (0.2,-1.9) rectangle (0.8,-2.6); \node[lbl,right] at (0.8,-2.25) {6\,kg};
\end{tikzpicture}}}
{\opt{$g$}{$g/2$}{$g/5$}{$g/10$}}

\QQ{A wire of length $L$, area of cross section $A$ is hanging from a fixed support. The length of the wire changes to $L_1$ when mass $M$ is suspended from its free end. The expression for Young's modulus is:}
{\opt{$\dfrac{MgL_1}{AL}$}{$\dfrac{Mg(L_1-L)}{AL}$}{$\dfrac{MgL}{AL_1}$}{$\dfrac{MgL}{A(L_1-L)}$}}

\QQ{The average thermal energy for a mono-atomic gas is: ($k_B$ is Boltzmann constant and $T$, absolute temperature)}
{\opt{$\dfrac12k_BT$}{$\dfrac32k_BT$}{$\dfrac52k_BT$}{$\dfrac72k_BT$}}

\QQ{Which of the following graph represents the variation of resistivity ($\rho$) with temperature ($T$) for copper?}
{\optfig{\gaxes{$T$}{$\rho$}{\draw (0,0.5) -- (1.7,1.4);}}
{\gaxes{$T$}{$\rho$}{\draw (0.1,1.6) .. controls (0.3,0.5) and (0.8,0.2) .. (1.5,0.1);}}
{\gaxes{$T$}{$\rho$}{\draw (0,0.1) .. controls (0.9,0.15) and (1.3,0.6) .. (1.4,1.6);}}
{\gaxes{$T$}{$\rho$}{\draw (0.1,1.5) .. controls (0.9,1.4) and (1.3,0.8) .. (1.6,0.1);}}}

\QQ{The color code of a resistance is given below:
\fig{\begin{tikzpicture}[scale=0.7]
\draw (-0.6,0.5) -- (0,0.5); \draw (5,0.5) -- (5.6,0.5); \draw (0,0) rectangle (5,1);
\foreach \x/\c in {0.9/Yellow,1.9/Violet,2.9/Brown,3.9/Gold} {\fill (\x-0.15,0) rectangle (\x+0.15,1); \node[lbl,below] at (\x,0) {\c};}
\end{tikzpicture}}
The values of resistance and tolerance, respectively, are:}
{\opt{470\,k$\Omega$, 5\%}{47\,k$\Omega$, 10\%}{4.7\,k$\Omega$, 5\%}{470\,$\Omega$, 5\%}}

\QQ{In Young's double slit experiment, if the separation between coherent sources is halved and the distance of the screen from the coherent sources is doubled, then the fringe width becomes:}
{\opt{double}{half}{four times}{one-fourth}}

\QQ{The capacitance of a parallel plate capacitor with air as medium is $6\unit{\mu F}$. With the introduction of a dielectric medium, the capacitance becomes $30\unit{\mu F}$. The permittivity of the medium is: ($\varepsilon_0 = 8.85\times10^{-12}\unit{C^2N^{-1}m^{-2}}$)}
{\opt{$0.44\times10^{-13}\unit{C^2N^{-1}m^{-2}}$}{$1.77\times10^{-12}\unit{C^2N^{-1}m^{-2}}$}{$0.44\times10^{-10}\unit{C^2N^{-1}m^{-2}}$}{$5.00\unit{C^2N^{-1}m^{-2}}$}}

\QQ{Dimensions of stress are:}
{\opt{$[\mathrm{MLT^{-2}}]$}{$[\mathrm{ML^2T^{-2}}]$}{$[\mathrm{ML^0T^{-2}}]$}{$[\mathrm{ML^{-1}T^{-2}}]$}}

\QQ{Assume that light of wavelength $600\unit{nm}$ is coming from a star. The limit of resolution of telescope whose objective has a diameter of $2\unit{m}$ is:}
{\opt{$3.66\times10^{-7}\unit{rad}$}{$1.83\times10^{-7}\unit{rad}$}{$7.32\times10^{-7}\unit{rad}$}{$6.00\times10^{-7}\unit{rad}$}}

\QQ{A series LCR circuit is connected to an ac voltage source. When L is removed from the circuit, the phase difference between current and voltage is $\dfrac\pi3$. If instead C is removed from the circuit, the phase difference is again $\dfrac\pi3$ between current and voltage. The power factor of the circuit is:}
{\opt{zero}{0.5}{1.0}{$-1.0$}}

\QQ{A short electric dipole has a dipole moment of $16\times10^{-9}\unit{C\,m}$. The electric potential due to the dipole at a point at a distance of $0.6\unit{m}$ from the centre of the dipole, situated on a line making an angle of $60^\circ$ with the dipole axis is: $\left(\dfrac1{4\pi\varepsilon_0} = 9\times10^9\unit{N\,m^2/C^2}\right)$}
{\opt{50\,V}{200\,V}{400\,V}{zero}}

\QQ{An iron rod of susceptibility 599 is subjected to a magnetising field of $1200\unit{A\,m^{-1}}$. The permeability of the material of the rod is: ($\mu_0 = 4\pi\times10^{-7}\unit{T\,m\,A^{-1}}$)}
{\opt{$2.4\pi\times10^{-4}\unit{T\,m\,A^{-1}}$}{$8.0\times10^{-5}\unit{T\,m\,A^{-1}}$}{$2.4\pi\times10^{-5}\unit{T\,m\,A^{-1}}$}{$2.4\pi\times10^{-7}\unit{T\,m\,A^{-1}}$}}

\QQ{A long solenoid of $50\unit{cm}$ length having 100 turns carries a current of $2.5\unit{A}$. The magnetic field at the centre of the solenoid is: ($\mu_0 = 4\pi\times10^{-7}\unit{T\,m\,A^{-1}}$)}
{\opt{$6.28\times10^{-4}\unit{T}$}{$3.14\times10^{-4}\unit{T}$}{$6.28\times10^{-5}\unit{T}$}{$3.14\times10^{-5}\unit{T}$}}

\QQ{A charged particle having drift velocity of $7.5\times10^{-4}\unit{m\,s^{-1}}$ in an electric field of $3\times10^{-10}\unit{V\,m^{-1}}$, has a mobility in $\mathrm{m^2V^{-1}s^{-1}}$ of:}
{\opt{$2.25\times10^{15}$}{$2.5\times10^6$}{$2.5\times10^{-6}$}{$2.25\times10^{-15}$}}

\QQ{The quantities of heat required to raise the temperature of two solid copper spheres of radii $r_1$ and $r_2$ ($r_1 = 1.5r_2$) through 1\,K are in the ratio:}
{\opt{$\dfrac{27}8$}{$\dfrac94$}{$\dfrac32$}{$\dfrac53$}}

\QQ{An electron is accelerated from rest through a potential difference of $V$ volt. If the de Broglie wavelength of the electron is $1.227\times10^{-2}\unit{nm}$, the potential difference is:}
{\opt{10\,V}{$10^2$\,V}{$10^3$\,V}{$10^4$\,V}}

\QQ{Taking into account of the significant figures, what is the value of $9.99\unit{m}-0.0099\unit{m}$?}
{\opt{9.9801\,m}{9.98\,m}{9.980\,m}{9.9\,m}}

\QQ{A ball is thrown vertically downward with a velocity of $20\unit{m/s}$ from the top of a tower. It hits the ground after some time with a velocity of $80\unit{m/s}$. The height of the tower is: ($g = 10\unit{m/s^2}$)}
{\opt{360\,m}{340\,m}{320\,m}{300\,m}}

\QQ{A capillary tube of radius $r$ is immersed in water and water rises in it to a height $h$. The mass of the water in the capillary is $5\unit{g}$. Another capillary tube of radius $2r$ is immersed in water. The mass of water that will rise in this tube is:}
{\opt{2.5\,g}{5.0\,g}{10.0\,g}{20.0\,g}}

\QQ{A resistance wire connected in the left gap of a metre bridge balances a $10\unit{\Omega}$ resistance in the right gap at a point which divides the bridge wire in the ratio 3:2. If the length of the resistance wire is $1.5\unit{m}$, then the length of $1\unit{\Omega}$ of the resistance wire is:}
{\opt{$1.0\times10^{-2}\unit{m}$}{$1.0\times10^{-1}\unit{m}$}{$1.5\times10^{-1}\unit{m}$}{$1.5\times10^{-2}\unit{m}$}}

\QQ{For the logic circuit shown, the truth table is:
\fig{\begin{circuitikz}[scale=0.7,transform shape]
\draw (0,1) node[not port](N1){}; \draw (0,-0.6) node[not port](N2){}; \draw (2.2,0.2) node[nor port](G){};
\draw (N1.in) -- ++(-0.5,0) node[left]{A}; \draw (N2.in) -- ++(-0.5,0) node[left]{B};
\draw (N1.out) -| (G.in 1); \draw (N2.out) -| (G.in 2); \draw (G.out) -- ++(0.4,0) node[right]{Y};
\end{circuitikz}}}
{\opt{\ttab{0&0&0}{0&1&0}{1&0&0}{1&1&1}}{\ttab{0&0&0}{0&1&1}{1&0&1}{1&1&1}}{\ttab{0&0&1}{0&1&1}{1&0&1}{1&1&0}}{\ttab{0&0&1}{0&1&0}{1&0&0}{1&1&0}}}

\end{multicols}
